% Version 0.29. One complete conditional case-II count chain.
% Pinned source: 7cbf7a0c32fca1badf05931879a37c3f8a2b03b4.
\section{Closing the case-II chain through the original count}
\label{sec:case-two-installation}

\status{The case-II seed is supplied by SX8, and UX12 constructs
the complete actual case-II promodular witness from that seed and
\texttt{PMExAlt}.  YX1 closes the resulting original-point
image-density statement in this case.  YX2 carries it through the
actual arithmetic datum, universal model, tangent comparison and
automorphic catalogue to the literal \texttt{CountII}.
The remaining input to this entire case-II chain is the actual
\texttt{PMExAlt} producer, with its narrower sufficient support
input analyzed in QX.  Case III, its distinct seed and its
component-reaching argument remain separate.}

% BEGIN YX1
\begin{theorem}[YX1: case-II image density from the single tower input]
\label{claim:YX1}
Assume the literal \texttt{PMExAlt}.  For every original
$D:\texttt{Raw}\,\Omega\,c$ with $\texttt{Ctx2}\,D$, and
every fixed $\iota:\overline{\mathbb Q}_\ell\simeq\mathbb C$,
\[
       D.\texttt{Fin0}\ \wedge\ D.\texttt{NoEnd}\ \wedge\
              D.\texttt{AllIrr}(\texttt{IsAut4Alt}\,\iota).
                                                               \tag{YX.1}
\]
In particular the one embedding required for the case-II part of
the source image-density statement can be chosen before all its
varying points and coefficient witnesses.
\end{theorem}
\begin{proof}
Fix $D$ and $\iota$ first.  The exact \texttt{Ctx2} includes
$D.\texttt{Good}$ and an original host, exponent, character and
\texttt{OfChar} witness with the prescribed parameters
$(n_{\mathrm{II}},1,1)$.  BX8 supplies the whole cohomology
record, and TW1 supplies the Ferrero--Washington input in the
scope used by UX.  SX8 supplies the literal \texttt{SeedIIAlt}
for every one of these towers and the same chosen $\iota$.
The assumed \texttt{PMExAlt} is its other arithmetic input.

Now take any original ambient $\mathfrak a$, any actual small
pair $(r,u)$ satisfying \texttt{KPair}, and suppose that the
representation in its ambient field is \texttt{Star}.
UX12 constructs an actual \texttt{ProModAlt} witness for this
pair.  This includes its tower, coefficient inclusion, original
\texttt{Space}, ideal, specialization into the original
valuation ring, conjugate representation, restrictions and local
characters.  Its coefficient threshold is chosen in the order
required by \texttt{PMExAlt}; none of these later choices changes
the already fixed $\iota$.
We have therefore supplied the universal star-point assertion
(IX.35) on this original $D$.

Here are its three uses in IX, to make the entire chain explicit.
For a characteristic-$\ell$ star pair, IX9 applies the finite
special-weight-fibre clause of its actual witness and obtains
constant residual characteristic polynomials.  The matrix
commutator obstruction excludes both irreducible full span and
the nonsplit \texttt{BlockExt} alternative.  This is the whole
$\operatorname{Con}(D)$ assertion.  In characteristic zero,
the same witness and IX10 give the precise weak Frobenius-root
ratio $\operatorname{PureW}$ on every star pair, including a
nonsplit extension constructed from a block diagonal point.
On full-span pairs IX10 also transports the actual automorphic
attachment and its \texttt{WitA} data to
\texttt{IsAut4Alt} at this same $\iota$.

IX8 now gives (YX.1).  Its proof includes the actual
characteristic-$\ell$ curve arising from failure of
\texttt{Fin0}, the integral nonsplit block construction, and
the weak-purity contradiction proving \texttt{NoEnd}.
The ordinary flags used there come from IX2 on the actual
points.  Thus the conclusion covers every original lift and
every original irreducible algebraic point in the required
quantifier order.  It uses no case-III witness.
\end{proof}
% END YX1

% BEGIN YX2
\begin{theorem}[YX2: the complete case-II count from the single tower input]
\label{claim:YX2}
With all the completed suppliers above, one has the literal
implication
\[
                         \boxed{\ \texttt{PMExAlt}
                                    \Longrightarrow\texttt{CountII}.\ }
                                                               \tag{YX.2}
\]
Precisely, for each primitive source character $\chi$ there is
$C_\chi\in\mathbb N$, chosen before
$h,t,\Omega,\ell,c,e,\psi$, such that every character in the
original Case II has finite strict Selmer length and
\[
        \frac{\ell_{c.O} V(\Omega,c,e,\psi)}{v_{c.O}(\ell)}
                    \log\ell\ \le (\ell h.N e)^{C_\chi}.
                                                               \tag{YX.3}
\]
Consequently \texttt{CountI}, \texttt{CountIII},
\texttt{CountV}, and \texttt{PMExAlt} suffice for the original
\texttt{PointHV}; the previously stated prime-pair hypotheses
then give the common cohomological-window consequences of
HV/TW/OB/VW.
\end{theorem}
\begin{proof}
Fix the tower input once.  CX8--CX9 supply the catalogue constants
$C,A$ before any host, coefficient ring or varying representation.
For a fixed primitive $\chi$, use the arithmetic absorption in
SP10 to choose the final exponent $C_\chi$ from these constants
and $\chi$.  The proof of that absorption uses only the source
host inequalities and the explicit level/exponent formulas below,
so this choice precedes all the remaining quantifiers.

Take an original host $h$ with $h.\chi=\chi$, an admissible
$t$, an odd $\ell$, the original $\Omega,c$, and
$\psi\in h.X(c,e)$ with $0<e<t$ in Case II.
First suppose that $O=c.O$ contains the finite dyadic root sets
required in AX5.  That construction gives an actual $D$ with
$D.\texttt{Good}$ and
\[
          \texttt{OfChar}(D,e,\psi,n',1,1),
                         \qquad n'=h.n_{\mathrm{II}}(\ell,\psi).
                                                               \tag{YX.4}
\]
Together with the data just fixed this is exactly
$\texttt{Ctx2}\,D$.  Choose one $\iota$ and apply YX1.
It supplies \texttt{Fin0}, \texttt{NoEnd} and the automorphy
of every irreducible actual point on this same datum.

AX7 gives the original common level and character exponent
\[
           L=\ell 2^{n'}h.N,
       \qquad k=(\ell-1)2^{n'}(2h.o).
                                                               \tag{YX.5}
\]
All three finite characters of $D$ satisfy
\texttt{KillsMu}(L) and have $k$th power one.
AX6, AX8 and AX9 give the actual continuity, ordinary weights,
level, dual twist and safety properties of every point.
Adding its automorphy from YX1 and its full matrix span gives
exactly the rank-four \texttt{AdmRAlt} record, for this one
$\iota$, the original $L$, and weights $(e+1,e,1,0)$.
This repeats the single pointwise use (AX.34) in AX10 and
completes G1 for Case II.

DX2, DX4 and DX8 give the original finiteness, dimensions and
local ring, so MX10 constructs the actual \texttt{Model}
over this same $O$.  Apply RX6 to each irreducible actual point,
using the automorphy supplied by YX1.  Its conclusion is
$D.\texttt{AllIrr}(D.\texttt{RigidAt})$, with all the original
local tangent conditions.  PX1--PX9 now supply the point,
conjugacy, orbit and tangent comparisons on this model.
With \texttt{Fin0} and \texttt{NoEnd} already available,
these are all of G2 used by SP10, and PX10 gives the actual
cotangent/trace inequality.  This supplies the geometric input
on the same original object as the arithmetic input.

For clarity, the resulting numerical step can be followed
without any additional image-density assumption.  Choose the
source prime $p=p_0(\ell h.N)\ge5$ and arithmetic Frobenius
$\sigma$ at the specified place.  Then $p\nmid L$ and
$p\notin D.S$.  The original diagonal representation has trace
\[
       \iota^{-1}(p^{e+1}z_1+p^e z_2+pz_3+1),\qquad z_j^k=1.
                                                               \tag{YX.6}
\]
CX9 gives a nonzero integer $N_*$ bounding the product of all
the relevant trace differences over any finite pairwise
nonisomorphic family.  Its full global labels, conjugacy
invariance and finite root multiplicities have already been
checked in CX; its enlarged auxiliary automorphic-label level
does not change the original Galois level $L$ in (YX.5).
PX10 and the same product inequality as in (JX.6) give
\[
 \frac{\ell_O D.\mathrm{Sha}}{v_O(\ell)}\log\ell
       \le -2\log|N_*|_\ell
       \le 2\log|N_*|
       \le 2Ck[L(e+1)]^A\log p.
                                                               \tag{YX.7}
\]
The finite-length assertion is part of the first step.
The factor two and the original valuation normalization
are retained.

The strict Selmer module in (YX.3) injects into the auxiliary
one used here by SP10's \texttt{shaEq} comparison.
At every newly admitted prime its negative-weight Frobenius
scalar is a root of unity times $p^{-e}$ and differs from $1$;
the scalar difference acts surjectively on $E/O$, giving the
needed local comparison.  SP10's arithmetic estimates absorb
(YX.7), with $p_1=1$ and the exact $n_{\mathrm{II}}$ in
(YX.4), into (YX.3).  Only the Case-II part of that argument
is used, so no case-III datum or seed enters the proof.

Finally let the initial $O$ be arbitrary.  Adjoin the two finite
dyadic root sets allowed in SP8 and work in a common finite
coefficient overring $O'$.  SP8 and RV1 preserve the kernel,
order, conductor, parity, case test and integer $n_{\mathrm{II}}$
of the original character, on the original algebraic closure.
The preceding construction applies over $O'$.  Those same
claims give equality of normalized strict Selmer lengths:
the length and $v(\ell)$ acquire the same ramification factor.
Returning to $O$ proves (YX.3) for its original coefficient
quantifier.  Neither $C_\chi$ nor any bound is changed by the
chosen extension.  This is exactly the source \texttt{CountShape}
with the Case-II predicate, hence the literal \texttt{CountII}.

For the last assertion combine this count with the three
stated counts and apply the original four-count assembly HV10.
The separately proved \texttt{CountVs}, \texttt{CountIV}, and
\texttt{CountZ} have their existing suppliers HV5, HV6 and CZ6.
The established window propagation retains its original field,
layer and prime-pair hypotheses.  In particular it makes no
claim that \texttt{CountIII} or the actual tower producer has
already been supplied.
\end{proof}
% END YX2
